% Deutsche Parallelfassung zu foundations/metric_realization.tex

\chapter{Metrische Realisierung}
\label{chap:metric-realization}

Nachdem die Krümmungsänderung einen überarbeiteten intrinsischen Zustand
erzeugt hat, ist die Trajektorie mathematisch bestimmt, aber noch nicht
räumlich messbar. Abstände zu Vermessungspunkten, World-Positionen und
Frame-Orientierungen erfordern einen metrischen Kontext. Metrische
Realisierung ist die Handlung, die diesen Kontext bereitstellt, ohne
Alignment-Identität oder konstruktive Geschichte zu verändern.

\section{Von intrinsischer Ordnung zu messbarer Geometrie}

Die intrinsische Konstruktion liefert die longitudinale Domäne, den
Krümmungsverlauf, geordnete Zustände und anwendbare Randbedingungen. Ein
metrischer Realization Context liefert den metrischen Raum, gegebenenfalls die
Referenzfläche, Einbettungskonventionen und weitere Informationen, die zur
Interpretation von Länge, Richtung, Krümmung und Position erforderlich sind.

\begin{definition}[Metrische Realisierung]
Metrische Realisierung ist die Interpretation einer intrinsischen
ingenieurtechnischen Beschreibung innerhalb eines bestimmten metrischen
Raums.
\end{definition}

Die Handlung wird geschrieben als
\[
\mathcal G_c:\mathcal X\rightarrow\mathcal X_{\mathrm{metric}},
\]
wobei der Index \(c\) den Realization Context ausdrücklich macht. Zwei
Kontexte können aus demselben intrinsischen Zustand verschiedene Koordinaten
oder metrische Folgen erzeugen und dabei dasselbe Engineering Object bewahren.

\paragraph{Ausführungsvertrag.}
Der Operator empfängt einen intrinsischen Zustand und einen Realization
Context; er erzeugt einen metrischen Zustand, in dem räumliche Größen bewertet
werden können. Er erfordert ein anwendbares metrisches Modell und hinreichende
Einbettungsdaten. Er bewahrt die Engineering Identity und trägt die Provenienz
sowohl des intrinsischen Zustands als auch des Kontexts. Er kann weder eine
physische Eisenbahn erzeugen, ihren Zustand beobachten, eine autoritative
Koordinatenquelle auswählen noch entscheiden, ob eine Abweichung akzeptabel
ist.

\section{Die fünf Rollen getrennt halten}

Im Beispiel der Krümmungsänderung sind die Rollen nun operativ getrennt:

\begin{description}
\item[Intrinsische Konstruktion] bestimmt, wie sich die bearbeitete Krümmung
  in Richtungswinkel und Trajektorie fortpflanzt.
\item[Metrischer Kontext] bestimmt, wie die Trajektorie messbare Distanz,
  Richtung, Krümmung und räumliche Position erhält.
\item[Koordinatenrepräsentation] hält das realisierte Ergebnis in einer
  gewählten Koordinatenbeschreibung fest; \(\mathcal C\) kann diese
  Beschreibung ändern, ohne die intrinsische Konstruktion zu wiederholen.
\item[Physische Realisierung] betrifft die gebaute oder instand gehaltene
  Eisenbahn und wird durch \(\mathcal R\), nicht durch \(\mathcal G_c\), zum
  Zielzustand in Beziehung gesetzt.
\item[Beobachtung] liefert quellengebundene Evidenz über einen Zustand,
  beispielsweise Vermessungsstichproben mit Provenienz und Unsicherheit.
\end{description}

Nach der Verarbeitung können diese Rollen dieselben numerischen Koordinaten
teilen. Numerische Übereinstimmung macht sie nicht austauschbar.

\section{Die Abweichung berechenbar machen}

Es sei \(x_{\mathrm{target}}\) der überarbeitete intrinsische Zielzustand. Der
metrische Zustand
\[
x_{\mathrm{target}}^{c}=\mathcal G_c(x_{\mathrm{target}})
\]
kann als CAD-Kurve oder abgetastete Polyline repräsentiert werden.
Vermessungsevidenz muss vor dem Vergleich in einem kompatiblen Kontext
interpretiert werden. Eine Koordinatentransformation kann zwei
World-Space-Beschreibungen zusammenführen; eine World-to-Track-Abbildung kann
eine gleisbezogene Interpretation bereitstellen. Keine der beiden Operationen
beweist gemeinsame Identität, gemeinsamen Zustand oder Quellenautorität.

Nachdem Kompatibilität und Anwendbarkeit begründet wurden, kann eine Bewertung
ein qualifiziertes Residuum erzeugen. Ein Wechsel von \(c\) kann dieses
Residuum verändern, weil sich Projektion, Flächenmodell oder metrische
Interpretation geändert hat. Diese Empfindlichkeit ist eine Eigenschaft der
Bewertungsqualifikation, kein Beleg dafür, dass das intrinsische Alignment
eine neue Identität erhalten hat.

\section{Warum das Ellipsoid folgt}

Eine lokal euklidische Realisierung kann für eine räumlich begrenzte
Ingenieuraufgabe ausreichen. Wo Projektionsverzerrung oder
Erdoberflächengeometrie wesentlich ist, muss der Realization Context
stattdessen die Referenzfläche und ihre metrische Struktur offenlegen. Die
Ellipsoidkapitel beginnen daher keine neue Geometriegeschichte. Sie
instanziieren den Vertrag der metrischen Realisierung für eine Fläche, auf der
Geradheit, Krümmung, Normalenrichtung und Frame-Transport
flächenbewusste Definitionen erfordern.
